
\def\PUDcodigo{ACE002}

\def\PUDcodacad{04507.1}

\def\PUDdisciplina{Álgebra Linear}

\def\PUDsemestre{S1}

\def\PUDhoras{80}

%NOVOS ---------------------------------------------------
\def\PUDhorasTeoria{80}
\def\PUDhorasPratica{0}

\def\PUDatuaisECA{---}

\def\PUDrecursos{Projetor multimídia; Quadro branco e pincel.}

%----------------------------------------------------------

\def\PUDcreditos{4}

\def\PUDprerequisito{---}

\def\PUDpreacad{---}

\def\PUDobjetivos{Compreender conceitos iniciais e resultados importantes da Álgebra linear, essenciais ao entendimento de outros conteúdos da matemática e da Engenharia, através dos seguintes objetivos específicos: Identificar matrizes, determinando a soma e o produto; Calcular determinante de uma matriz, matriz adjunta e matriz inversa; Resolver sistemas de equações lineares, relacionando com as matrizes; Identificar transformações lineares, determinando o núcleo e a imagem; Definir operadores lineares, calculando autovalores e autovetores de um operador linear, identificando o polinômio característico de uma matriz e o polinômio diagonalizável.}

\def\PUDementa{Matrizes; Determinantes; Sistemas de Equações Lineares; Espaços Vetoriais; Transformações Lineares; Autovalores e autovetores.}

\def\PUDconteudo{ & Matrizes:  Matrizes.  Tipos especiais de matrizes.  Operações com matrizes.  Aplicações. 
\\ 
 &  Determinantes:        Conceitos preliminares.  Determinantes.  Desenvolvimento de Laplace.  Matriz adjunta e matriz inversa.  Regra de Cramer.  Aplicações.  
\\ 
 &  Sistemas lineares:       Sistemas de equações lineares.  Sistemas lineares e matrizes.  Operações elementares com linhas ou colunas de uma matriz.  Matriz na forma escada.  Resolução de sistemas de equações lineares.  Inversão de matrizes.  Aplicações.  
\\ 
 & Transformações lineares:        Espaço vetorial.  Funções vetoriais.  Transformações lineares.  Núcleo de uma transformação linear.  Imagem de uma transformação linear.  Matriz de uma transformação linear.  Aplicações. 
\\ 
 &  Autovalores e Autovetores:        Operadores lineares.  Autovalores e autovetores de um operador linear.  Polinômio característico.  Diagonalização de operadores.  Aplicações. }

\def\PUDmetodo{Aulas expositórias, em que busca-se sempre mostrar aplicações da disciplina, de forma a servir como estopim para a curiosidade do aluno, motivando-o a estudar e aprofundar-se no mundo dessa subárea da Matemática. Além disso, busca-se, sempre que possível, envolver a monitoria na metodologia de ensino, de tal forma que melhore o aprendizado do aluno, amenizando assim suas dificuldades e os ajudando a progredir no curso. 
%Aulas expositórias que podem ser teóricas e/ou práticas, onde as práticas no laboratório serão marcadas no decorrer da disciplina.
}

\def\PUDavalia{A avaliação será feita com aplicação de provas teóricas, além da possibilidade de inclusão de trabalhos, seminários e listas de exercícios no decorrer da disciplina.
%A avaliação será feita com aplicação de provas teóricas e/ou práticas, além da possibilidade de inclusão de trabalhos e seminários no decorrer da disciplina.
}

\def\PUDbibbasica{ &  \href{www.biblioteca.ifce.edu.br/index.asp?codigo_sophia=63261}{Anton}, Howard.  Álgebra linear com aplicações.  10ª ed.  Porto Alegre: Bookman, 2012.  768p.  ISBN: 9788540701694. 
\\ 
 &  \href{www.biblioteca.ifce.edu.br/index.asp?codigo_sophia=44606}{Steinbruch}, Alfredo.  Introdução à álgebra linear.  São Paulo: Pearson Makron Books, 2013.  245p.  ISBN: 9780074609446. 
\\ 
 &  \href{www.biblioteca.ifce.edu.br/index.asp?codigo_sophia=23859}{Boldrini}, José Luiz.  Álgebra linear.  3º ed.  São Paulo: Harbra, 1986.  411 p.  ISBN: 8529402022. }

\def\PUDbibcomplementar{ &  \href{www.biblioteca.ifce.edu.br/index.asp?codigo_sophia=63180}{Lipschutz}, Seymour.  Álgebra Linear.  4º ed.  Porto Alegre, RS : Bookman, 2011.  ISBN: 9788577808335. 
\\ 
 &   \href{www.biblioteca.ifce.edu.br/index.asp?codigo_sophia=41150}{Kolman}, Bernard.  Introdução a Álgebra Linear com Aplicações.  8º ed. Rio de Janeiro, RJ : LTC, 2014.  664p.  ISBN: 9788521614784. 
\\ 
 &  \href{www.biblioteca.ifce.edu.br/index.asp?codigo_sophia=40456}{Strang}, Gilbert.  Álgebra linear e suas aplicações.  São Paulo, SP: Cengage Learning, 2013.  ISBN: 9788522107445.
\\
 &  \href{www.biblioteca.ifce.edu.br/index.asp?codigo_sophia=58225}{Fernandes}, Daniela Barude.  Álgebra Linear. E-book.  Pearson.  146p.  ISBN: 9788543009568.
\\
 &  \href{www.biblioteca.ifce.edu.br/index.asp?codigo_sophia=59843}{Franco}, Neide Maria Bertoldi. Álgebra linear. E-book. Pearson. 376p.  ISBN: 9788543019154. 
}

\def\PUDautor{Samuel Vieira Dias}

\def\PUDdata{2013-04-15}

\def\PUDrevisao{1}

\def\PUDrevisor{Samuel Vieira Dias}

\def\PUDdatarevisao{2017-05-02}

\PUDcorpo
